\scanpage{28}
2) On a aussi $\mathrm{Hom}_V(A,B)=\mathrm{Hom}_U(U\otimes_V A,B)$

Ce qui nous ramène à la \inserted{situation} \uline{duale} de \uline{celle} \deleted{que} envisagée dans les généralités précédentes. \deleted{\unsure{Mais}}

[On voit donc que les objets \deleted{\unsure{$S$}} $\sim$-projectifs sont les facteurs directs d'objets $U\otimes_V A$.]

Ici les hom. canoniques \deleted{\unsure{sont}}

\[
\begin{aligned}
&\deleted{U\otimes_V A\to A}\\
e_A^*:\quad & A\longrightarrow U\otimes_V A \qquad (V\text{-linéaire, donné par }a\unsure{-|\to}1\otimes a)
\end{aligned}
\]

\[
i_A^{+}:\quad U\otimes_V A-A \qquad (U\text{-linéaire surjectif par }u\otimes_V a-\unsure{ua|au})
\]

3) Soit $\mathcal C$ la catégorie des faisceaux de $\underline O$-modules sur $X$, et $\widetilde{\mathcal C}$ la catégorie produit $\prod_{x\in X}\mathcal C^{\underline O_x}$

On a la \unsure{solution|situation} du pb. universel grâce au foncteur de Godement $F-C^0(F)$. \quad ---

% paper watermark: ARGENTINE TYPEWRITER (vertical; reverse of sheet shows through)
% archive sheet number: 27
